
\section{Twenty-Fifth Sunday after Trinity: The Coming of the Son of Man}

\begin{quote}

Matthew 24:35-44. Heaven and earth shall pass away, but My words shall by no means pass away. But of that day and hour no one knows, not even the angels of heaven, but the Father alone. But as the days of Noah were, so shall also the coming of the Son of Man be. For as they were in the days before the Flood, eating and drinking, marrying and giving in marriage, until the day that Noah entered into the ark, and they took no heed, until the Flood came and swept them all away, so shall also the coming of the Son of Man be. Then shall two be in the field; the one shall be taken, and the other left. Two women shall be grinding at the mill; the one shall be taken, and the other left. Watch therefore! For you do not know in what hour your Lord comes. But know this, that if the master of the house had known in what watch of the night the thief would come, he would have watched, and would not have let his house be broken into. Therefore be you also ready! For the Son of Man comes in an hour when you do not think.

\end{quote}

\bigskip


“Creation has been subjected to vanity,” says Paul. “All flesh is grass, and all its glory as the flower of the grass; the grass withers, and its flower falls away.” All life in the world comes to an end; it bears the seed of death within itself, and vanity has set its stamp upon it. Like a garment the heavens shall be rolled together and changed, and the earth shall perish.

This is as certain as it is appointed unto men once to die. But here God’s Gospel and the Christian faith do not come to a halt. If there were nothing more to preach, one might well say: that sorrow can wait until its time. But now there is more, infinitely more. For God’s Gospel does not preach death, but life. Therefore this is not the main thing in our text, that heaven and earth shall pass away; but this is the main thing: “My words shall by no means pass away.” It is not the end of the world, but the coming of the Son of Man, that is the main thing for the Christian faith. Therefore this cry shall sound for the awakening of all who sleep: The Lord comes! Behold, the Lord God comes with might, and His arm rules; behold, His reward is with Him, and His recompense before Him. The Son of Man comes; and He comes in the hour when you do not think; He comes to judge the living and the dead.

But it is written of the people in Noah’s days that “they ate and drank, married and gave in marriage, until the day that Noah entered into the ark, and they took no heed, until the Flood came and swept them all away.”

So shall the coming of the Son of Man be. “They took no heed”: that was the dreadful frivolity with which Noah’s contemporaries answered the Lord’s solemn warning through Noah’s mouth. And then the Flood came, and then there was crying and misery and fear and flight—and no deliverance any more. Death-struggle and death-agony for a little while, and then the stillness of death, when all was over.

“They took no heed”: that was the frivolity with which the Jews answered the warning cry of John the Baptist and of Jesus and of Jesus’ disciples: “Repent, for the kingdom of heaven is at hand.” They took no heed; they defied God’s voice, and they despised Jesus’ tears over stubborn Jerusalem. And the judgment came, and the city was taken, and the temple was burned, and there was not left one stone upon another; and the unhappy people was driven into exile and bondage, sold into the hand of its enemies.

“They took no heed.” And alas, it appears that the same way is being taken yet once again. For again it goes as in the days of Noah: “They ate and drank, married and gave in marriage.” These words do not describe any exceptional viciousness; but they denote complete frivolity, false security, worldliness, and godlessness. This condition can indeed be joined with the very highest human civilization, and even with a very respectable morality; but not with the fear of God and faith in Jesus Christ. Therefore the world would so gladly have Christ’s Gospel out of the world. The world wants peace and quiet for its sensual enjoyments; it will not be disturbed by thoughts of death and thoughts of judgment and the earnestness of eternity. To enjoy life, to enjoy sensuality and vanity, that is the matter; away with the Gospel of the coming of the Son of Man! Thus it goes as in the days of Noah: “They ate and drank, married and gave in marriage.”

Poor, blinded child of man! What will you say when the end of these things comes? For the end comes, the Lord comes, and what shall then take place? Then comes the judgment, the separation, the decision, the eternal decision: “The one shall be taken, and the other shall be left.” Then it goes according to the inflexible rule: “The Lord knows those who are His.” His elect He shall gather from the four corners of the world, and not one of them shall be forgotten or passed over. But the wicked and all stumbling blocks shall be given over to the fire that is not quenched, and the worm that does not die.

How will it go with you on that day? Do you know the Lord, and does the Lord know you? Or are you still the world’s friend and God’s enemy? Little will the friendship of the world help you on that day, when the world passes away with its desire. As Belshazzar’s feast came to an end in terror in the night when Babylon fell, so the world’s joy will abruptly come to an end when the day of the Lord comes as a thief in the night upon all those who dwell on the face of the whole earth.

Therefore the loving voice of the blessed Gospel sounds to you: “Awake, you who sleep! Awake, awake!” Do you think that you will live upon the earth forever, or do you think that the earth will stand to all eternity, since you order your life as though this were your abiding place? “Fool, this night your soul is required of you.” What will you do when death and judgment come, you who have used your whole life and all your powers in the service of the world and of sin?

Awake, while it is the time of grace, and remember Jesus Christ, who is raised from the dead. He is the one who comes again; He is the one who shall hold judgment. What have you done with Him who went into death in order to save you out of the world’s corruption and open for you the way into God’s blessedness? If you have rejected love, then love will cast you off. Therefore awake and seek grace in the time of grace; behold, now is the accepted time; behold, now is the day of salvation.

But also to God’s children, and quite especially to them, this word sounds today: Watch therefore! There is nothing more necessary in the present time of this world than watchfulness. And yet it is almost forgotten among us. The watchmen slumber upon the walls, and the inhabitants of the city dream sweet dreams and forget the danger. Or is there not a mist of worldliness that enshrouds us all, and a cradle-song of the world that lulls us into sweet sleep? Where are those watchmen’s voices that cry aloud and spare not, but declare to the Lord’s people their sins? Where are those watchful Christians who stand in the whole armor of God, ready to fight against sin in all its forms? Worldliness grows strong in the midst of God’s congregation, and love grows cold in many because the desire of the world grows strong. There are many to whom the Lord sends this word: “But I have this against you, that you have left your first love. Remember therefore from where you have fallen, and repent, and do the first works; but if not, then I come to you soon and will remove your lampstand from its place, if you do not repent.”

Watch therefore! Let the heart turn away from the world and its vain joys and sorrows, and be lifted up toward the Lord, the soul’s eternal Bridegroom, who soon comes to lead the Bride home. Only that one is prepared who waits for the Bridegroom; only that one is prepared who has longing for heaven in the soul and rejoices to depart from the world and its vanity; only that one is prepared for whom Jesus is the first and the last, beginning and end.

Watch therefore! “For the Son of Man comes in the hour when you do not think.”

